\section{Síntesis y ampliación}

\subsection{Conclusiones principales}

\begin{enumerate}
\item AI for Math no consiste en resolver ejercicios de matemáticas en serie, sino en un sistema completo que va del problema y la intuición a la formalización y la verificación.
\item Los proof assistants como Lean/Coq le dan a la IA una retroalimentación estricta, lo que convierte a las matemáticas en un dominio de alto valor, adecuado para el entrenamiento y la evaluación.
\item La intuición matemática y la verificación formal no se oponen; un sistema sólido necesita contar a la vez con free attention y bounded attention.
\item Un Neo Lab como Axiom solo se sostiene si confluyen una red de matemáticos, un sistema de formalización, el entrenamiento de modelos, financiamiento y capacidad de cómputo, y un objetivo de producto.
\item La meta de AI for Math no es sustituir a los matemáticos, sino construir junto con ellos demostraciones, bibliotecas de teoremas y líneas de investigación.
\end{enumerate}

\subsection{Preguntas abiertas}

\begin{itemize}
\item ¿Podrá la IA generar demostraciones con verdadero valor estético y poder explicativo, y no solo demostraciones largas y correctas?
\item ¿La formalización en Lean se convertirá en la interfaz principal de la investigación matemática, o solo servirá a algunas áreas?
\item ¿Cómo puede escalarse la retroalimentación de los matemáticos para que se forme un volante de datos de AI for Math?
\item ¿El producto de un Neo Lab como Axiom tomará la forma de un proof assistant, de una plataforma de investigación o de un sistema de colaboración entre matemáticos?
\end{itemize}

\subsection{Lecturas adicionales}

\begin{itemize}
\item Lean theorem prover y mathlib: para entender la infraestructura de las matemáticas formalizadas.
\item Las discusiones públicas de Terence Tao, Kevin Buzzard y otros sobre formalization: para entender por qué la comunidad matemática le da importancia a Lean.
\item Artículos sobre AI theorem proving, autoformalization y proof search: para entender las rutas técnicas de AI for Math.
\item Las notas de EP139 y EP138: para situar AI for Math dentro del panorama más amplio de los agentes y los laboratorios de modelos.
\end{itemize}

\begin{importantbox}{Valoración final}
La dificultad de AI for Math reside en que exige al modelo, al mismo tiempo, entender matemáticas, escribir en un lenguaje formal, someterse a la verificación automática y colaborar con la intuición de los matemáticos. Es el campo de prueba decisivo para que la IA pase de ser una «herramienta de productividad» a ser una «compañera en el descubrimiento científico».
\end{importantbox}

\end{document}
